% Translator-owned textual collation and independent mathematical verification.
\head{第15节前半：参数族、非空性公理及维数}

\noindent\textbf{实际范围。}
本模块从法文手稿物理90页末的第15节标题开始，完整覆盖91、92页，并至93页“与（2.1）至（2.11）相应的结果”一句结束；英文整理稿对应59页标题至61页同一句。下一句关于（2.12）的条件及本节后半的微分论证另列报告。本轮实际显示查看了法90--93和英59--61的页面图像，另放大读取法91开头、92页首公式及手添词、92页末与93页首；法3--5页另作第1节有限型范围的旁证。原作者的删改、英编脚注68--70、本轮修正和补证分别标明。

当前旧稿\path{continuations/v22_section15_inherited.tex}及其原有V-H-L048--069校注完整保留，不回写为本轮结果。正文采用十五条V23-SET001--015新注，英编68--70另外以英译编者注保存。本轮不是第15节的首次翻译，而是对已有译文的法英回校、补足英文漏译及有证据的订正。

\subsection*{逐段对应}
\begin{longtable}{@{}p{.14\linewidth}p{.18\linewidth}p{.57\linewidth}@{}}
\toprule 标记 & 法／英物理页 & 原段及处理\\\midrule\endhead
SET-P01 & 90末--91上／59下 & 公理化研究计划；恢复正式成稿最初几节及“下文”，保留第4、6节原号。\\
SET-P02 & 91中／59末 & 有限表现交换图$(D)$；恢复下划线对象，补充箭头名供核查。\\
SET-P03 & 91下--92首／60上 & 参数解释、剩余域扩张、$T$值点、$B_T$订为$P_T$；保留英编注68。\\
SET-P04 & 92上／60中 & $\mathcal X=X\times_P\mathcal P$及交换图、提炼公理的任务；恢复\textit{dégager}，保留英编注69。\\
SET-P05 & 92中／60中下 & 平坦、光滑及几何不可约纤维；恢复手添\textit{ici et plus tard}，保留英编注70。\\
SET-P06 & 92下／60下 & 逆像维数、像维数不足时为空；显明有限型、非空纤维及参数范围。\\
SET-P07 & 92末--93上／60末 & 非空性作为新公理；法稿$\ne$及$f(Z)$均正确，英文误写。\\
SET-P08 & 93上中／61首 & 一般闭集的占优与维数；给出逐分支公式及混合维反例。\\
SET-P09 & 93中／61上 & （2.1）至（2.11）的推广概括；恢复范围号，限制局部性质的正则性前提。\\
\bottomrule
\end{longtable}

\audit{V23-A01}{法稿恢复、英文异文及原稿删改}{能恢复的文字不再留作缺文，错误归属逐项区分}
法90--91说“主要想到第2至第8节的结果”，并明确建议在正式成稿的最初几节即采用公理化观点。英文59页漏了这一编排建议。法91的\textit{ci-dessous}指下文，英文却写“above”；这影响微分条件的指向，已据法稿恢复。关于Serre、Zariski的叙述保留第一人称及建议口吻，不能改成译者已采取的行动，或把“我有这样的印象”写成已证定理。

首段法91确实写第4、6节，本节末法94则写第3、6节。正文保留两处原号，报告指出差异；没有足够依据把首段的4静默改为3。这里涉及Bertini--Zariski推广的研究安排，不能仅凭编号差别就宣称原作者的数学命题矛盾。

法91以带下划线的$P$表示关联预概形，法92又区别带下划线的$X,Z$。英文转写丢掉这些区分，才形成“$P\to P$”“$X=X_P\times P$”等难读式子。本版用$\mathcal P,\mathcal X,\mathcal Z$统一转排，是记号恢复，不是新添对象。$(D)$及两个交换方形据原页保留；箭头$q,p,q_X,p_X$是明示的编辑命名，不声称原图已有这些字母。

英编稿关于\textit{grassmannienne}“是形容词！”的括注与脚注68保留为英编意见。法文\textit{une grassmannienne induite}中的用法是名词。法91的$s$位于$\xi$“下方”也清楚可读，英文“over”倒置了结构态射方向。相反，法92的目标确实印为$B_T$，不是仅英文误读；由于$B$没有定义且拉回构造的目标只能是$P_T$，此项按法稿的记号错误订正。

法92的\textit{dégager}完整可辨，英文脚注69的词义试探不应继续造成译文缺词。该页在\textit{envisagées}后手添的\textit{ici et plus tard}也可读，意为“这里及后文”；英文所谓“长段字迹不清”不能继续作为法文缺文处理。对应英编脚注70仍保留，以便知道前一版为什么曾标疑。

法92的$\dim f(Z)<m$与法93的$\mathcal Z_\eta\ne\varnothing$均正确，英文分别误作$f(X)$与等号。因此这两项属于英文转写订正，不能算作原作者定理的两个错误。法93的“2.1.\ à 2.11.”表示连续范围，而英文“2.1 and 2.11”只留下两端。法92末划去的维数加强句可辨，移入正文脚注说明；法93三行密集覆打删除不全可辨，本版没有重构其内容或从中猜造额外假设。

\audit{V23-A02}{参数族、基变换及非空纤维}{形式构造成立，交换性与笛卡尔性分开}
记图$(D)$中的态射为
\[
q:\mathcal P\longrightarrow P,\qquad
p:\mathcal P\longrightarrow G,
\qquad j=(q,p):\mathcal P\longrightarrow P\times_SG.
\]
给定$S$态射$\xi:T\to G$，定义
\[
\mathcal P_\xi=\mathcal P\times_{G,\xi}T.
\]
交换性恰好给出一个$T$态射$\mathcal P_\xi\to P_T$。如果$j$为闭浸入，该态射也是闭浸入，因为它是$j$的基变换。若$T=\Spec k(\xi)$，则$P_T$正是$P_s\times_{k(s)}k(\xi)$。这里无需$T$为域的谱，无需$\xi$是截面，也无需$(D)$本身为笛卡尔图。

对任意$f:X\to P$，$\mathcal X=X\times_P\mathcal P$及其态射到$X$、$G$都由纤维积给出。$q_X:\mathcal X\to X$是$q$沿$f$的基变换，所以光滑性、几何不可约纤维及其维数$N-m$都保留。这个形式构造没有有限型限制；但若要求新方形的每条箭头都有限表现，就必须另外确保$X\to S$有相应的有限表现条件。不能把“同类型的图式”当作这一有限性结论的证明。

本节“几何不可约纤维”按非空不可约的约定解释。因此$q$的每个纤维非空，$q$满射；再由光滑性，$q$是普遍开映射。开性的准确一般版本可核对\href{https://stacks.math.columbia.edu/tag/01UA}{Stacks Project，Tag 01UA}。几何不可约性及光滑性在域扩张后保持，使每个纤维成为几何整的光滑预概形；这里“几何整”是关于$q$的纤维，并不是关于$G$的额外断言。

参数$m$表示所取截面的余维范围，取$m\geq0$；纤维非空且有维数$N-m$又强制$m\leq N$。这两个边界值得明写：若允许$m=-1$，取$P=\Spec k$、$G=\mathbb A^1_k$、$\mathcal P=\mathbb A^2_k$及$N=1$，令$q$为结构态射、$p$为常值$0$。此时光滑纤维维数为$N-m=2$，而“每个$m$维闭不可约子集”的公理没有任何对象可检查。取$X=Z=P$，虽有$\dim f(Z)=0\geq m$，泛截面却为空。故$m\geq0$不能在应用公理时被忘记，正文不将这一负余维范围纳入本节。

对于不可约的闭子预概形$W\subset P$，$q_W:\mathcal P\times_PW\to W$为开满射且纤维不可约，故总空间不可约。直接证明如下：两个非空开集的像是$W$中的非空开集，二者相交；在交中的一个纤维上，两个原开集又给出非空开子集，不可约性使它们相交。更一般地，逆像使$W$与总空间的不可约分支一一对应，可核对\href{https://stacks.math.columbia.edu/tag/037A}{Stacks Project，Tag 037A}。这一论证将在非空性公理的推广中使用。

\audit{V23-A03}{维数公式的有限型范围}{逆像公式正确，任意非有限型外推不成立}
从这里起$S=\Spec k$。图$(D)$的有限表现条件使$P,G,\mathcal P$都在$k$上有限型。第1节法稿物理5页已明确把$X$放在有限型范围；第15节在形式构造处只写“一个态射$f$”，容易掩盖后续维数推理所使用的有限性。因此正文在开始维数论证时将$X/k$有限型明写出来。因为$P/k$有限型，这也包含第1节所说的$f$有限型的通常情形；本版不就法5页右边缘残损处的$P$或$S$字形作新的确定判断。

令$Z$是$X$的闭子集，赋予约化闭子预概形结构；写成有限个不可约分支$Z_i$的并。对应的$\mathcal Z_i$不可约、在$Z_i$上占优且一般纤维维数$N-m$。利用有限型域代数的维数等于函数域超越次数，得到
\[
\dim\mathcal Z_i
 =\dim Z_i+(N-m).
\]
这里的代数事实见\href{https://stacks.math.columbia.edu/tag/00P0}{Stacks Project，Tag 00P0}；同一计算也可由光滑态射的局部维数公式得到，见\href{https://stacks.math.columbia.edu/tag/02FW}{Tag 02FW}。对有限个分支取最大值即得
\[
\dim\mathcal Z=\dim Z+N-m.
\]
所以原稿在这一处允许$Z$不纯维并无错误。若$\dim Z<m$，总空间维数小于$\dim G=N$，不可能占优。若只知道$\dim\overline{f(Z)}<m$，则先对$W=\overline{f(Z)}\subset P$使用同一计算；$q^{-1}(W)\to G$不占优，而$\mathcal Z\to G$经由这个空间，故其泛纤维仍为空。

这项维数推理确实不能脱离有限型范围。取
\[
P=G=\mathbb A^1_k,\qquad
\mathcal P=\Delta(\mathbb A^1_k),\qquad N=m=1,
\]
两个投影均为同构。令$X=\Spec k(t)$，并用$k[u]\hookrightarrow k(t)$、$u\mapsto t$定义$f$。作为预概形，$X$只有一个点而$\dim X=0$，但它映到$G$的泛点，且
\[
\mathcal X_\eta=\Spec k(t)\ne\varnothing.
\]
这反驳了在完全任意$X$上使用“$d<m$故不能占优”。它没有反驳第1节已经明确的有限型版本，也不反驳对任意$f$成立的纤维积构造。

$G$的非约化结构不影响本项纯维数判断：取约化后底层空间、维数和泛点剩余域不变，有限型维数计算可在约化的不可约对象之间进行。但是这不能用于证明与平坦性或正则性有关的结论，见V23-A06。

\audit{V23-A04}{从$m$维闭集公理到一般非空性}{补足像与闭包之间的论证，公理不由初始条件自动推出}
采用非空性公理
\[
\text{对每个闭不可约 }W\subset P,\quad
\dim W=m\ \Longrightarrow\ (q^{-1}W)_\eta\ne\varnothing.
\tag{A}
\]
给闭集赋予约化结构只用于说明纤维积，其非空性及维数都与这个选择相容。公理(A)对于这里的光滑、平坦条件是一个额外要求。例如取$P=G=\mathbb A^1_k$、$\mathcal P=\mathbb A^1_k$、$q$为恒等、$p$为常值$0$、$N=m=1$，则初始条件都成立，甚至$j$是到$P\times G$的闭浸入，但$W=P$的泛截面为空。因此原稿将(A)称作“基本公理”是有内容的，不能改译为前述假设的推论。

现设$Z\subset X$不可约且$\dim f(Z)\geq m$，令$T=\overline{f(Z)}$。由于有限型，$f(Z)$是$T$中的稠密可构造集，故包含一个非空开集$U\subset T$。所用的是\href{https://stacks.math.columbia.edu/tag/054K}{Stacks Project，Tag 054K，Chevalley定理}及\href{https://stacks.math.columbia.edu/tag/053X}{Tag 053X，稠密可构造集包含非空开集}。仅知道$f(Z)$稠密，不足以省略这一步。

由于$0\leq m\leq\dim T$，可在$T$中选择一个$m$维闭不可约子集$W$与$U$相交。具体地，在$U$中取非空仿射开集，作Noether正规化，然后在维数为$\dim T$的仿射空间中取一个$m$维坐标子空间；有限占优映射的逆像有一个支配该子空间的不可约分支。其在$T$中的闭包具有维数$m$，且仍与$U$相交。因为$T$闭于$P$，$W$也闭于$P$，故可以应用(A)。

由V23-A02，$\mathcal W=q^{-1}(W)$不可约；由(A)，它在$G$上占优。$q^{-1}(W\cap U)$是$\mathcal W$的非空开集，因此含其泛点，其到$G$的态射仍占优。另一方面，$f^{-1}(U)\cap Z\to U$满射，任意基变换后仍满射：对每个点，非空的剩余域纤维在进一步域扩张后仍非空。沿$q^{-1}(W\cap U)\to U$作基变换，得到一个满射，其源是$\mathcal Z$的子空间。因此$\mathcal Z\to G$也占优，所需的泛截面非空性成立。

可约闭集$Z$的情形对其有限个不可约分支应用上述论证即可。若$\dim f(Z)\geq m$，至少有一个分支的像维数达到$m$，故至少有一个非空泛截面。结合V23-A03得到准确的判据
\[
\mathcal Z_\eta=\varnothing
\quad\Longleftrightarrow\quad
\dim f(Z)<m.
\]
空集可单独处理，或采用$\dim\varnothing=-\infty$的约定。这里证明的是本固定域及本参数族下的非空性判据；原稿关于以后各项公理在任意基变换下稳定的总括语句，不由本段补证自动获得一个未陈述的一般定理。

\audit{V23-A05}{泛纤维维数与可约分支}{原一概$\dim Z-m$的断言错误，改为实际存活分支的最大值}
若$Z$不可约且$\dim f(Z)\geq m$，则$\mathcal Z$不可约并在$G$上占优。设$\zeta$为$\mathcal Z$的泛点；在取约化后，函数域塔给出
\[
\begin{aligned}
\dim\mathcal Z_\eta
 &=\trdeg_{k(\eta)}k(\zeta)\\
 &=\trdeg_k k(\zeta)-\trdeg_k k(\eta)\\
 &=\dim\mathcal Z-\dim G
 =\dim Z-m.
\end{aligned}
\]
泛纤维不可约，是总空间在$G$的泛点上的局部化再取同一底层空间；计算并不要求$k$代数闭。若$G$非约化，泛点剩余域仍与$G_{\mathrm{red}}$的函数域相同，纯维数等式可在约化对象上验证。

如果$Z$可约，记其不可约分支为$Z_i$。$\dim f(Z_i)<m$的分支在泛截面中消失；其余分支的泛截面不可约，维数为$\dim Z_i-m$。它们相互不能包含：每个非空泛截面的泛点也是相应总空间的泛点，并映到$Z_i$的泛点；若它包含于另一分支的逆像，便强迫$Z_i$包含于另一$Z_j$。于是它们恰是$\mathcal Z_\eta$的不可约分支，且
\[
\dim\mathcal Z_\eta
 =\max_{\{i:\,\dim f(Z_i)\geq m\}}(\dim Z_i-m).
\tag{*}
\]
这里已假定至少一个分支存活。如果所有分支都消失，则泛截面为空。原来的$\dim Z-m$等式，当且仅当至少有一个最大维分支存活时成立；要求每一个分支的像维数都至少$m$是一个较强但方便的充分条件。

反例可以完全处在原稿的有限型、闭关联族范围。取
\[
P=G=\mathbb P^1_k,\qquad
\mathcal P=\Delta(\mathbb P^1_k),\qquad N=m=1,
\]
并令
\[
X=\mathbb P^1_k\amalg\mathbb A^2_k.
\]
$f$在第一分支上是恒等，在第二分支上恒等于一个$k$有理点$p\in P$。$q$、$p$作为关联族的两个投影均为同构，因而光滑且纤维几何不可约、维数$0=N-m$；图中的全部态射有限表现，$P,G$在$k$上平坦，$j$是闭浸入。$P$的唯一一维闭不可约子集是$P$本身，它的泛截面非空，故(A)成立。

但对$Z=X$，有$\dim Z=2$、$\dim f(Z)=1\geq m$；第二分支的常值像不会遇到$G$的泛点，第一分支贡献一个点，因而
\[
\mathcal Z_\eta=\Spec k(\mathbb P^1),\qquad
\dim\mathcal Z_\eta=0\ne2-1.
\]
这个反例精确否定原稿对任意闭集的维数等式，并不否定占优性、非空性判据或不可约情形的等式。正文同时保留一般闭集的叙述与公式(*)，没有把可约情形仅藏在校注中。

法93把“占优”的源写成$\mathcal Z_\eta$。从拓扑上说，非空泛纤维到$G$的像就是泛点，故该态射确实占优；问题不在于“占优”一词本身为假，而在于随后减去$\dim G$的维数式所用的对象是有限型的总空间$\mathcal Z$。正文明确总空间，避免把通常不是有限型$k$预概形的泛纤维拿来再减一次$\dim G$。

\audit{V23-A06}{推广局部性质时的非约化参数障碍}{维数论证不能替代正则态射前提}
原稿在得到非空性与维数性质后，笼统写它们允许展开（2.1）至（2.11）的相应结果。这里必须把纯拓扑、维数结论与依赖局部环的结论分开。前者已经在上文核验；后者原证明使用泛截面到$X$的正则态射性质，不能仅从“$G$几何不可约”得到。

取$A=k[\epsilon]/(\epsilon^2)$，并令
\[
P=G=\mathcal P=X=\Spec A,
\qquad p=q=f=\id,
\qquad N=m=0.
\]
所有箭头均有限表现，$P,G$在域$k$上平坦；$G$几何不可约且维数$0$，$q$为同构，光滑且每个几何纤维是一个点。$P$唯一的非空闭不可约子集维数为$0$，其泛截面非空，故(A)也成立。$G$只有一个点$\eta$，但
\[
\mathscr O_{G,\eta}=A\ne k(\eta)=k.
\]
于是$\mathcal X_\eta=\Spec k$，而$\mathcal X_\eta\to X$对应商映射$A\to k$，不是平坦态射。例如把单射$(\epsilon)\hookrightarrow A$张量上$k$后，非零的
$(\epsilon)\otimes_Ak\simeq k$
映成零，故张量函子不保持这一单射。$A$的Krull维数为$0$而极大理想需要一个生成元，所以$A$不正则；$k$是域，正则。因此（2.11）中“截面在一点正则当且仅当$X$在像点正则”的局部结论在这套一般公理下失败。约化性的相应等价也失败。不能把这个现象当作英文下划线丢失造成的纯记号问题。

一个足够的补充条件是$G$在泛点处约化。由于$G$不可约且有限型，此时
\[
\mathscr O_{G,\eta}=k(\eta).
\]
在任一含$\eta$的仿射开集中，$\Spec k(\eta)\to G$就是在泛素理想处的局部化，所以平坦；它的非空几何纤维是域上的一个几何正则点。这使它成为正则态射。正则态射的定义及光滑态射的正则性可核对\href{https://stacks.math.columbia.edu/tag/07R6}{Stacks Project，Tag 07R6}。本处沿$p_X$的基变换仍为局部化型态射，故
\[
\mathcal X_\eta\longrightarrow\mathcal X
\]
平坦，其非空几何纤维是点。$q_X$又光滑；所以复合
\[
h:\mathcal X_\eta\longrightarrow X
\]
平坦，并且它在$X$的一点上的纤维是$q_X$的光滑纤维的局部化。这些局部化是诺特的，在任意剩余域扩张后仍为正则环，故$h$为正则态射。对$h$上的一对点，其局部环同态是平坦局部同态，故忠实平坦。这恢复了原来局部性质传递论证所需要的正则性前提。

上述条件只在本固定域的泛点处要求约化，比把整个$G$预先改成光滑或几何整弱；原稿没有这项条件，正文把它明示为足够的补充，而不冒称原文已有。若还要在任意扩域后统一使用这项补充，须另核基变换后的泛点条件，不能把一般约化与几何约化混淆。本轮没有逐条重新审定（2.1）至（2.11）的全部历史证明，也没有利用这一补充抹去其中已有的个别勘误和有限性限制。后半节在另加$P$光滑之后研究泛截面的光滑性，是另一项论证，见本版后半报告。
